\BourbakiExerciseGroup

\begin{enumerate}
\def\labelenumi{\arabic{enumi})}
\tightlist
\item
  Let A, B be two subsets of a topological space X such that \(A \supset B\) . Show that:
\end{enumerate}

\begin{enumerate}
\def\labelenumi{\alph{enumi})}
\item
  The interior of B with respect to X is contained in the interior of B with respect to the subspace A of X. Give an example where these two sets are distinct.
\item
  The frontier of B with respect to A is contained in the intersection of A with the frontier of B with respect to X. Give an example where these two sets are distinct.
\end{enumerate}

\begin{enumerate}
\def\labelenumi{\arabic{enumi})}
\setcounter{enumi}{1}
\tightlist
\item
  Let A and B be any two subsets of a topological space X.
\end{enumerate}

\begin{enumerate}
\def\labelenumi{\alph{enumi})}
\item
  Show that the intersection of A and the interior of B with respect to X is contained in the interior of \(B \cap A\) with respect to A. Give an example where these two sets are distinct.
\item
  Show that the intersection of A and the closure of B with respect to X contains the closure of B ∩ A with respect to A. Give an example where the two sets are distinct.
\end{enumerate}

\begin{enumerate}
\def\labelenumi{\arabic{enumi})}
\setcounter{enumi}{2}
\item
  A subspace \(\mathbf{A}\) of a topological space \(\mathbf{X}\) is discrete if and only if every point of \(\mathbf{A}\) is isolated.
\item
  Let Y, Z be two subspaces of a topological space X such that \(X = Y \cup Z\) and such that \(\overline{A} \cap B = \overline{B} \cap A = \varnothing\) , where \(A = Y \cap CZ\) and \(B = Z \cap CY\) .
\end{enumerate}

\begin{enumerate}
\def\labelenumi{\alph{enumi})}
\item
  Show that if M is any subset of X then the closure of M in X is the union of the closure of \(M \cap Y\) in Y and the closure of \(M \cap Z\) in Z. Deduce that if \(M \cap Y\) is closed (resp. open) in Y and \(M \cap Z\) is closed (resp. open) in Z, then M is closed (resp. open) in X.
\item
  Let \(f\) be a mapping of \(\mathbf{X}\) into a topological space \(\mathbf{X}'\) . Show that if \(f|_{\mathbf{Y}}\) and \(f|_{\mathbf{Z}}\) are continuous then so is \(f\) .
\end{enumerate}

\begin{enumerate}
\def\labelenumi{\arabic{enumi})}
\setcounter{enumi}{4}
\item
  Let Y, Z be two subspaces of a topological space X such that \(X = Y \cup Z\) . Show that if a subset M of \(Y \cap Z\) is open (resp. closed) in both Y and Z, then M is open (resp. closed) in X.
\item
  Let A be a locally closed subspace of a topological space X. Show that the set of open subsets U of X such that \(A \subset U\) and A is closed in U has a greatest element, namely the complement in X of the frontier of A with respect to \(\overline{A}\) .
\item
  Show by examples that the union of two locally closed subsets and the complement of a locally closed subset need not be locally closed.
\item
  Let A be a subset of an ordered set X. Show that the topology induced on A by the right (resp. left) topology of X (§ 1, Exercise 2) is the right (resp. left) topology of A with respect to the ordering of A induced by that of X.
\item
  Let \(\mathbf{A}\) be a subset of a linearly ordered set \(\mathbf{X}\) .
\end{enumerate}

\begin{enumerate}
\def\labelenumi{\alph{enumi})}
\item
  Show that the topology induced on A by the topology \(\mathcal{G}_{-}(\mathbf{X})\) {[}resp. \(\mathcal{G}_{+}(\mathbf{X}),\mathcal{G}_{0}(\mathbf{X})]\) (§ 2, Exercise 5) is finer than the topology \(\mathcal{G}_{-}(\mathbf{A})\) {[}resp. \(\mathcal{G}_{+}(\mathbf{A}),\mathcal{G}_{0}(\mathbf{A})]\) , the ordering of A being that induced by the ordering of X. If A is an interval of X, then \(\mathcal{G}_{-}(\mathbf{A})\) {[}resp. \(\mathcal{G}_{+}(\mathbf{A}),\mathcal{G}_{0}(\mathbf{A})]\) is the topology induced on A by \(\mathcal{G}_{-}(\mathbf{X})\) {[}resp. \(\mathcal{G}_{+}(\mathbf{X}),\mathcal{G}_{0}(\mathbf{X})]\) .
\item
  Take X to be the product \(Q \times Z\) , ordered by the lexicographic ordering (Set Theory, Chapter III, § 2, no. 6) and take A to be the subset \(Q \times \{0\}\) . Show that the topology \(\mathcal{G}_{-}(A)\) {[}resp. \(\mathcal{G}_{+}(A)\) , \(\mathcal{G}_{0}(A)\) {]} is strictly coarser than the topology induced on A by \(\mathcal{G}_{-}(X)\) {[}resp. \(\mathcal{G}_{+}(X)\) , \(\mathcal{G}_{0}(X)\) {]}.
\end{enumerate}

\begin{enumerate}
\def\labelenumi{\arabic{enumi})}
\setcounter{enumi}{9}
\item
  Let A be a non-empty subspace of a non-empty topological space X. Show that the topology induced on \(\mathfrak{P}_{0}(\mathbf{A})\) by the topology \(\mathcal{G}_{\Omega}\) (resp. \(\mathcal{G}_{\Phi}\) ) on \(\mathfrak{P}_{0}(\mathbf{X})\) (§ 2, Exercise 7) is the topology \(\mathcal{G}_{\Omega}\) (resp. \(\mathcal{G}_{\Phi}\) ) on \(\mathfrak{P}_{0}(\mathbf{A})\) .
\item
  Let X be a topological space, \(\mathcal{G}\) the topology of X, A a dense subspace of X. Show that the set of topologies on X which are finer than \(\mathcal{G}\) , in which A is dense in X, and which induce the same topology as \(\mathcal{G}\) on A, has at least one maximal element; such a maximal element is called an A-maximal topology. Show that a topology \(\mathcal{G}_0\) on X is A-maximal if and only if every subset M of X, such that A∩M is dense in M (with respect to \(\mathcal{G}_0\) ) and open in A, is open in the topology \(\mathcal{G}_0\) , and that the subspace \(\mathbb{C}\) A is then discrete in the topology induced by \(\mathcal{G}_0\) , and is closed in X. Show also that if \(\mathcal{G}_0\) is A-maximal and if the topology induced by \(\mathcal{G}_0\) on A is quasi-maximal (§ 2, Exercise 6) then \(\mathcal{G}_0\) is quasi-maximal.
\end{enumerate}

* 12) Let X be the subspace {[}0, 1{]} of the real line R, and let S be the equivalence relation on X whose equivalence classes are \{0, 1\} and the sets \{x\} for 0 \textless{} x \textless{} 1. Show that X has no continuous section with respect to S. *

\begin{enumerate}
\def\labelenumi{\arabic{enumi})}
\setcounter{enumi}{12}
\item
  Let X be a topological space, and R and S two equivalence relations on X such that R implies S. Show that if there is a continuous section of X with respect to R and a continuous section of X/R with respect to S/R, then there is a continuous section of X with respect to S.
\item
  Let X be a topological space which is the sum of two subspaces \(X_{1}\) and \(X_{2}\) , and let X/R be the space obtained by pasting together \(X_{1}\) and \(X_{2}\) along a subspace \(A_{1}\) of \(X_{1}\) and a subspace \(A_{2}\) of \(X_{2}\) by means of a homeomorphism h of \(A_{1}\) onto \(A_{2}\) . Show that the canonical images of \(X_{1}\) and \(X_{2}\) in X/R are homeomorphic to \(X_{1}\) and \(X_{2}\) respectively.
\end{enumerate}

* 15) Consider the following three subspaces of \(\mathbf{R}: \mathbf{X}_1 = ] - \mathbf{1}, \mathbf{1}[, \mathbf{X}_2 = ] - \mathbf{2}, -\mathbf{1}[, \mathbf{X}_3 = ]\mathbf{1}, \mathbf{2}[.\) Let \(\mathbf{X}\) be their union; \(\mathbf{X}\) is then the sum of these three subspaces. Let \(A_{2}\) (resp. \(A_{3}\) ) be the set of irrational numbers contained in \(X_{2}\) (resp. \(X_{3}\) ); let \(B_{1}\) (resp. \(B_{2}\) ) be the set of rational numbers contained in {]}-1, 0{[} (resp. \(X_{2}\) ); and let \(C_{1}\) (resp. \(C_{3}\) ) be the set of rational numbers contained in {]}0, 1{[} (resp. \(X_{3}\) ). Let X/R be the space obtained by pasting together the \(X_{i}\) (i) along \(A_{2}\) and \(A_{3}\) by means of the homeomorphism \(x \to x + 3\) ; (ii) along \(B_{1}\) and \(B_{2}\) by means of the homeomorphism \(x \to x - 1\) ; (iii) along \(C_{1}\) and \(C_{3}\) by means of the homeomorphism \(x \to x + 1\) . Show that the canonical image of \(X_{1}\) in X/R is not homeomorphic to \(X_{1}\) . If \(Y = X_{1} \cup B_{2} \cup C_{3}\) is the saturation of \(X_{1}\) with respect to R, show that the quotient space \(Y/R_{Y}\) is not homeomorphic to the canonical image of Y in X/R. *

\begin{enumerate}
\def\labelenumi{\arabic{enumi})}
\setcounter{enumi}{15}
\tightlist
\item
  Let X be a non-empty topological space and R an equivalence relation on X. Show that if X/R is regarded as a subset of \(\mathfrak{P}_{0}(X)\) , then the quotient topology on X/R is coarser than the topologies induced on X/R by the topologies \(C_{\Omega}\) and \(C_{\Phi}\) on \(\mathfrak{P}_{0}(X)\) (§ 2, Exercise 7).
\end{enumerate}

